\section{无穷级数}

\subsection{无穷级数的收敛性}

\begin{definition}
    对于一个无穷项数列 $\{a_n\}$，我们有部分和数列 $S_n = \sum_{i=1}^{n} a_i$。那么我们称
    $$
    \sum_{i=1}^{\infty} a_i = \lim_{n \to \infty} S_n
    $$
    是一个无穷级数。当 $\lim_{n \to \infty} S_n$ 存在时称级数收敛，否则称级数发散。
\end{definition}

无穷级数具有一系列性质：

\begin{theorem}
    如果级数 $\sum\limits_{n = 1}^{\infty} a_n$ 收敛，那么 $\lim\limits_{n \to \infty} a_n = 0$。

    或等价描述为，若 $\lim\limits_{n \to \infty} a_n \neq 0$，那么 $\sum\limits_{n = 1}^{\infty} a_n$ 发散。
\end{theorem}

\begin{theorem}[线性性]
    如果级数 $\sum\limits_{n = 1}^{\infty} a_n$ 和 $\sum\limits_{n = 1}^{\infty} b_n$ 都收敛，那么 $\sum\limits_{n = 1}^{\infty} (\lambda a_n + \mu b_n)$ 也收敛，且
    $\sum\limits_{n = 1}^{\infty} (\lambda a_n + \mu b_n) = \lambda \sum\limits_{n = 1}^{\infty} a_n + \mu \sum\limits_{n = 1}^{\infty} b_n$。
\end{theorem}

\begin{theorem}
    若 $\sum\limits_{n = 1}^{\infty}$ 收敛，那么任意加括号之后组成的新级数也收敛。
\end{theorem}

\begin{theorem}
    若给级数加括号后收敛，并且括号中各项符号相同，那么原级数也收敛于同一个值。

    \begin{proof}
        我们考虑加括号后的数列
        $$
        b_k = \sum_{i = n_{k - 1} + 1}^{n_k} a_i
        $$
        设 $\{b_k\}$ 的部分和为 $A_1, A_2, \dots, A_k, \dots$，原级数法部分和是 $S_1, S_2, \dots, S_k, \dots$, $\lim\limits_{n \to \infty} A_n = S$。那么当 $n_{k - 1} + 1 \le n \le n_k$ 时，有：
        $$
        A_{k - 1} \le S_n \le A_k \lor A_{k - 1} \ge S_n \ge A_k
        $$
        于是得证。
    \end{proof}
\end{theorem}

利用 Cauchy 收敛法则，我们可以得到无穷级数的 Cauchy 收敛定理：

\begin{theorem}[Cauchy 收敛定理]
    $\sum\limits_{i = 1}^{\infty} a_i$ 收敛等价于：
    $$
    \forall\,\varepsilon > 0, \exists\,N \in \mathbb{N}, \forall\,m \ge n \ge N, \abs{\sum_{i = n}^{m} a_i} < \varepsilon
    $$
\end{theorem}

\subsection{正项级数}

\begin{definition}[正项级数]
    对于级数 $\sum_{n = 0}^{\infty} a_n$，若 $a_n \ge 0$，那么称其为正项级数。
\end{definition}

显然，正项级数的收敛等价于它有界。

\subsection{正项级数的比较判别法}

\begin{theorem}[比较判别法]
    对于正项级数 $\{a_n\}, \{b_n\}$。设 $\exists\,N: \forall\,n > N: 0 \le a_n \le b_n$，那么

    \begin{itemize}
        \item 若 $\sum\limits_{n = 1}^{\infty} b_n$ 收敛，那么 $\sum\limits_{n = 1}^{\infty} a_n$ 也收敛；
        \item 若 $\sum\limits_{n = 1}^{\infty} b_n$ 发散，那么 $\sum\limits_{n = 1}^{\infty} a_n$ 也发散。
    \end{itemize}
\end{theorem}

\begin{theorem}[比较判别法的极限形式]
    对于正项级数 $\{a_n\}, \{b_n\}$。设 $l = \lim\limits_{n \to \infty} \frac{a_n}{b_n}$，那么：
    \begin{itemize}        
        \item 若 $l = 0$，那么当 $\sum_{n = 0}^{\infty} b_n$ 收敛时 $\sum_{n = 0}^{\infty} a_n$ 收敛；
        \item 若 $l > 0$，那么当 $\sum_{n = 0}^{\infty} b_n$ 与 $\sum_{n = 0}^{\infty} a_n$ 同敛散；
        \item 若 $l = +\infty$，那么当 $\sum_{n = 0}^{\infty} b_n$ 发散时 $\sum_{n = 0}^{\infty} a_n$ 发散。
    \end{itemize}
\end{theorem}
